\documentclass[sigconf,nonacm]{acmart}

\hypersetup{
  pdfsubject={Computational Cognitive Modeling final project proposal},
}

\begin{document}

\title{Beyond the Dyad: Linguistic Style Transfer\texorpdfstring{\\}{ }Across Conversational Partners}

% ===== 小组成员：替换为实际姓名，人数不同时增删即可 =====
% 如需加单位，在每个 \author 后加 \affiliation{\institution{...}\country{...}}（acmart 要求 country 必填）
\author{Eric Wu}
\author{Hao Su}
\author{Luis Zhang}

\maketitle
% 单页文档：在第一栏就调用 \balance，让最后一页两栏底部对齐
\balance

People can pick up ways of speaking from a conversational partner. This project asks whether those habits carry over when they write to someone else. \citet{Kronm_ller_2026} distinguished adaptation to individual speakers from adaptation to the wider conversation in a reference task. In a controlled production task, \citet{Ostrand_2019} found that participants' syntax shifted with the forms they heard across partners. We extend this question to several dimensions of style in naturally interleaved conversations. We predict that after receiving messages from A, a player P will write to B in a style closer to A's, even after accounting for P's earlier usage and the messages B has sent to P.

The public \textit{Diplomacy} corpus \citep{peskov-etal-2020-takes} contains 17,289 private messages from 12 seven-player games. We replay each game in message order, predicting the style of each outgoing message from the sender, recipient, and prior conversation. Once observed, the message becomes part of the recipient's history. To measure style, we mask place names and move orders, then extract a standardized vector of features such as pronoun use, hedges, contractions, and punctuation. Paired rewrites will test whether the vector captures changes in expression when content stays fixed and remains stable when only the content changes \citep{Halvani_2021,wegmann-etal-2022-author}. Our primary analysis considers the whole vector, with individual features examined separately.

A hierarchical Gaussian model tracks what a player learns from incoming messages, drawing on Hawkins et al.'s account of shared and partner-specific expectations and Xu and Futrell's model of how linguistic preferences update with experience \citep{Hawkins_2023,xu2024hierarchicalbayesianmodelsyntactic}. To predict a message, each model blends the player's early writing style with the style estimate it uses for that recipient. A weight fitted to observed messages determines the size of the shift. The no-adaptation baseline fixes that weight at zero. If A's style changes while B's history stays fixed, a partner-specific learner predicts a shift only in messages to A. Full pooling predicts comparable shifts to both partners under matched conditions. Partial pooling allows A's influence to carry over to B, with a smaller shift when P has already received substantial input from B.

In our pilot, stylistic associations were clearer within a partnership, where the input also tended to be more recent. We will compare inputs at similar distances in the message sequence while accounting for game phase, message length, and reply context. Messages unavailable to P provide a negative control for shared circumstances. In simulations, we vary A's input with B's history fixed and redistribute identical inputs across partners to separate exposure from source. On games held out one at a time, higher predictive log scores for models with shared learning and shifts toward A in messages to B would support cross-partner adaptation. The observational results will help identify which predictions to test under controlled exposure.

\bibliographystyle{ACM-Reference-Format}
\bibliography{references}

\end{document}
